\chapter{RESULTS}

\section{Characteristics of the study population}

The final analytic sample comprised 4,102 women aged 15--29 years from the Vietnam MICS6 survey (2020--2021). Table~\ref{tab:descriptive} presents the survey-weighted distribution of sociodemographic characteristics.

\begin{table}[htbp]
\caption{Weighted sociodemographic characteristics of the study population (N~=~4,102)}
\centering
\small
\begin{tabular}{|p{4.5cm}|p{4.0cm}|c|}
\hline
\textbf{Characteristic} & \textbf{Category} & \textbf{\%} \\
\hline
\multirow{3}{*}{Age group} & 15--19 & 30.4 \\
 & 20--24 & 29.7 \\
 & 25--29 & 39.9 \\
\hline
\multirow{2}{*}{Residence} & Urban & 37.4 \\
 & Rural & 62.6 \\
\hline
\multirow{2}{*}{Ethnicity} & Kinh/Hoa & 85.0 \\
 & Ethnic minority & 15.0 \\
\hline
\multirow{3}{*}{Marital status} & Currently married/union & 46.4 \\
 & Formerly married/union & 3.3 \\
 & Never married/union & 50.3 \\
\hline
\multirow{3}{*}{Education} & Primary or less & 5.2 \\
 & Lower secondary & 22.1 \\
 & Upper secondary+ & 72.7 \\
\hline
\multirow{5}{*}{Wealth quintile} & Q1 (poorest) & 19.2 \\
 & Q2 & 21.2 \\
 & Q3 & 21.9 \\
 & Q4 & 20.1 \\
 & Q5 (richest) & 17.6 \\
\hline
\multirow{2}{*}{Health insurance} & Insured & 86.3 \\
 & Uninsured & 13.7 \\
\hline
\multirow{2}{*}{Mass media exposure} & Yes & 90.9 \\
 & No & 9.1 \\
\hline
\multirow{2}{*}{ICT exposure} & Yes & 97.9 \\
 & No & 2.1 \\
\hline
\multirow{2}{*}{Decision autonomy} & High & 90.2 \\
 & Low & 9.8 \\
\hline
\multirow{2}{*}{Ever had sexual intercourse} & Yes & 53.1 \\
 & No & 46.9 \\
\hline
\hline
\multicolumn{3}{|l|}{\textbf{Outcome variables}} \\
\hline
HPV vaccine awareness & Yes & 62.4 \\
HPV vaccination & Yes & 7.5 \\
\hline
\end{tabular}
\label{tab:descriptive}
\end{table}

The majority of women resided in rural areas (62.6\%), belonged to the Kinh/Hoa ethnic group (85.0\%), and half were never married (50.3\%). Most had attained upper secondary education or higher (72.7\%). Health insurance coverage was high (86.3\%), as was mass media exposure (90.9\%) and ICT access (97.9\%).

Regarding the primary outcomes, 62.4\% (95\% CI: 59.8--64.9\%) of women reported having heard of HPV vaccination, while only 7.5\% (95\% CI: 6.3--8.8\%) had received the vaccine---revealing a substantial awareness--vaccination gap.

\section{Bivariate analysis of HPV awareness}

Table~\ref{tab:awareness_bivariate} presents the weighted prevalence of HPV vaccine awareness by sociodemographic characteristics.

\begin{table}[htbp]
\caption{HPV vaccine awareness by sociodemographic characteristics (N~=~4,102)}
\centering
\small
\begin{tabular}{|p{3.8cm}|c|c|c|}
\hline
\textbf{Characteristic} & \textbf{Aware (\%)} & \textbf{95\% CI} & \textbf{p-value} \\
\hline
\multicolumn{4}{|l|}{\textit{Age group}} \\
\hline
15--19 & 48.3 & 44.2--52.4 & \multirow{3}{*}{$<$0.001} \\
20--24 & 66.2 & 61.9--70.1 & \\
25--29 & 70.3 & 67.0--73.4 & \\
\hline
\multicolumn{4}{|l|}{\textit{Residence}} \\
\hline
Urban & 73.1 & 69.0--76.8 & \multirow{2}{*}{$<$0.001} \\
Rural & 56.0 & 52.6--59.3 & \\
\hline
\multicolumn{4}{|l|}{\textit{Ethnicity}} \\
\hline
Kinh/Hoa & 67.0 & 64.2--69.8 & \multirow{2}{*}{$<$0.001} \\
Ethnic minority & 36.1 & 31.7--40.6 & \\
\hline
\multicolumn{4}{|l|}{\textit{Education}} \\
\hline
Primary or less & 22.9 & 16.0--31.6 & \multirow{3}{*}{$<$0.001} \\
Lower secondary & 49.2 & 44.3--54.1 & \\
Upper secondary+ & 69.2 & 66.4--71.9 & \\
\hline
\multicolumn{4}{|l|}{\textit{Wealth quintile}} \\
\hline
Q1 (poorest) & 33.3 & 29.0--37.9 & \multirow{5}{*}{$<$0.001} \\
Q2 & 58.5 & 53.6--63.3 & \\
Q3 & 64.4 & 59.3--69.1 & \\
Q4 & 73.5 & 68.8--77.8 & \\
Q5 (richest) & 83.6 & 79.5--87.0 & \\
\hline
\multicolumn{4}{|l|}{\textit{Mass media exposure}} \\
\hline
No & 39.1 & 32.8--45.7 & \multirow{2}{*}{$<$0.001} \\
Yes & 64.7 & 62.1--67.3 & \\
\hline
\multicolumn{4}{|l|}{\textit{Health insurance}} \\
\hline
Insured & 63.0 & 60.3--65.6 & \multirow{2}{*}{0.233} \\
Uninsured & 58.7 & 51.7--65.4 & \\
\hline
\multicolumn{4}{|l|}{\textit{Decision autonomy}} \\
\hline
High & 65.2 & 62.4--67.8 & \multirow{2}{*}{$<$0.001} \\
Low & 45.2 & 38.5--52.0 & \\
\hline
\end{tabular}
\label{tab:awareness_bivariate}
\end{table}

A clear socioeconomic gradient emerged: awareness increased from 33.3\% in the poorest quintile to 83.6\% in the richest ($p < 0.001$). Awareness rose sharply with education (22.9\% primary vs. 69.2\% upper secondary+). Furthermore, awareness was substantially higher among urban (73.1\%) compared to rural (56.0\%) women. Ethnic minority women had awareness of just 36.1\% compared to 67.0\% among Kinh/Hoa women. Notably, health insurance was not significantly associated with awareness ($p = 0.233$).

\begin{figure}[htbp]
\centering
\includegraphics[width=0.90\textwidth]{Figures/fig1_awareness_vaccination_by_ses.png}
\caption{HPV vaccine awareness and vaccination prevalence by wealth quintile among women aged 15--29, Vietnam MICS6 (2020--2021). A steep socioeconomic gradient is evident for both outcomes.}
\label{fig:awareness_vacc_ses}
\end{figure}

\begin{figure}[htbp]
\centering
\includegraphics[width=0.90\textwidth]{Figures/fig3_awareness_vaccination_by_age.png}
\caption{HPV vaccine awareness and vaccination by age group. Awareness peaks at age 25--29, while vaccination is highest among women aged 20--24.}
\label{fig:awareness_vacc_age}
\end{figure}

\section{Bivariate analysis of HPV vaccination}

Table~\ref{tab:vaccination_bivariate} summarizes vaccination prevalence by sociodemographic factors.

\begin{table}[htbp]
\caption{HPV vaccination status by sociodemographic characteristics (N~=~4,102)}
\centering
\small
\begin{tabular}{|p{3.8cm}|c|c|c|}
\hline
\textbf{Characteristic} & \textbf{Vaccinated (\%)} & \textbf{95\% CI} & \textbf{p-value} \\
\hline
\multicolumn{4}{|l|}{\textit{Age group}} \\
\hline
15--19 & 5.6 & 3.9--8.0 & \multirow{3}{*}{0.015} \\
20--24 & 10.0 & 7.7--12.8 & \\
25--29 & 7.0 & 5.5--8.9 & \\
\hline
\multicolumn{4}{|l|}{\textit{Residence}} \\
\hline
Urban & 11.4 & 9.1--14.1 & \multirow{2}{*}{$<$0.001} \\
Rural & 5.1 & 3.9--6.7 & \\
\hline
\multicolumn{4}{|l|}{\textit{Ethnicity}} \\
\hline
Kinh/Hoa & 8.5 & 7.2--10.1 & \multirow{2}{*}{$<$0.001} \\
Ethnic minority & 1.6 & 0.8--3.4 & \\
\hline
\multicolumn{4}{|l|}{\textit{Education}} \\
\hline
Primary or less & 0.8 & 0.1--4.0 & \multirow{3}{*}{$<$0.001} \\
Lower secondary & 2.4 & 1.3--4.3 & \\
Upper secondary+ & 9.5 & 8.0--11.3 & \\
\hline
\multicolumn{4}{|l|}{\textit{Wealth quintile}} \\
\hline
Q1 (poorest) & 0.7 & 0.3--1.7 & \multirow{5}{*}{$<$0.001} \\
Q2 & 6.0 & 4.1--8.9 & \\
Q3 & 5.9 & 4.1--8.5 & \\
Q4 & 9.2 & 6.7--12.4 & \\
Q5 (richest) & 16.6 & 12.7--21.3 & \\
\hline
\end{tabular}
\label{tab:vaccination_bivariate}
\end{table}

The wealth gradient in vaccination was extremely steep: only 0.7\% of women in the poorest quintile had been vaccinated, compared to 16.6\% in the richest---a 24-fold difference. This transition from a 2.5-fold disparity in awareness to a 24-fold disparity in vaccination demonstrates that out-of-pocket costs and health system access are the true rate-limiting steps. Vaccination was highest among women aged 20--24 (10.0\%), likely reflecting opportunistic catch-up vaccination.

\begin{figure}[htbp]
\centering
\includegraphics[width=0.88\textwidth]{Figures/fig4_awareness_vaccination_by_area.png}
\caption{HPV awareness and vaccination by urban/rural residence. Urban women show substantially higher rates for both outcomes.}
\label{fig:awareness_vacc_area}
\end{figure}

Figure~\ref{fig:awareness_vacc_area} graphically illustrates this geographic divide. Urban residence significantly facilitates both awareness and actual vaccine receipt.

\begin{figure}[htbp]
\centering
\includegraphics[width=0.88\textwidth]{Figures/fig5_awareness_vaccination_by_edu.png}
\caption{HPV awareness and vaccination by education level. A steep positive gradient is evident, with upper secondary+ women showing the highest rates.}
\label{fig:awareness_vacc_edu}
\end{figure}

As depicted in Figure~\ref{fig:awareness_vacc_edu}, educational attainment operates as a powerful structural determinant. Women with upper secondary education or higher demonstrate substantially greater vaccine uptake compared to those with primary education or less.

\begin{figure}[htbp]
\centering
\includegraphics[width=0.88\textwidth]{Figures/fig6_awareness_vaccination_by_ethnicity.png}
\caption{HPV awareness and vaccination by ethnicity. Ethnic minority women have substantially lower awareness and near-zero vaccination coverage.}
\label{fig:awareness_vacc_ethnicity}
\end{figure}

Figure~\ref{fig:awareness_vacc_ethnicity} highlights the marginalization of ethnic minority women, with awareness at 36.1\% and vaccination at just 1.6\%.

\section{Socioeconomic inequality: concentration index}

\subsection{Overall concentration indices}

Table~\ref{tab:ci_overall} presents the overall concentration indices for each outcome.

\begin{table}[htbp]
\caption{Wagstaff concentration indices for HPV awareness, vaccination, and awareness--vaccination gap}
\centering
\begin{tabular}{|p{4.0cm}|c|c|c|c|}
\hline
\textbf{Outcome} & \textbf{N} & \textbf{CI} & \textbf{SE} & \textbf{95\% CI} \\
\hline
Awareness & 4,102 & 0.149 & 0.009 & 0.131--0.166 \\
Vaccination & 4,102 & 0.389 & 0.047 & 0.296--0.482 \\
Gap & 2,133 & $-$0.035 & 0.006 & $-$0.047--$-$0.023 \\
\hline
\end{tabular}
\label{tab:ci_overall}

\vspace{0.2cm}
\footnotesize{CI = Wagstaff concentration index, ranked by household wealth score. Survey-weighted estimation using \texttt{conindex} in Stata. Positive values indicate pro-rich concentration; negative values indicate pro-poor concentration.}
\end{table}

All three concentration indices were statistically significant ($p < 0.001$). The positive CIs for awareness (0.149) and vaccination (0.389) indicate \textbf{pro-rich inequality}---both outcomes are concentrated among wealthier women. The magnitude of inequality was markedly greater for vaccination than for awareness, suggesting that wealth-related barriers operate more strongly at the point of vaccine uptake than at the point of information exposure.

The negative CI for the awareness--vaccination gap ($-$0.035) indicates that the gap is concentrated among poorer women who, despite being aware, face additional barriers to actually receiving the vaccine.

\subsection{Concentration curves}

\begin{figure}[htbp]
\centering
\includegraphics[width=0.88\textwidth]{Figures/fig7_concentration_curve.png}
\caption{Concentration curves for HPV awareness and vaccination, ranked by household wealth score. Both curves lie below the line of equality, confirming pro-rich concentration. The vaccination curve departs further from equality, reflecting greater socioeconomic inequality in vaccine uptake.}
\label{fig:concentration_curves}
\end{figure}

As shown in Figure~\ref{fig:concentration_curves}, both concentration curves lie below the diagonal line of equality, visually confirming that awareness and vaccination are disproportionately concentrated among wealthier women. The vaccination curve lies further from the equality line than the awareness curve, consistent with the larger CI value for vaccination (0.389 vs. 0.149).

\section{Multivariable Poisson regression}

Table~\ref{tab:poisson} presents the adjusted prevalence ratios from multivariable survey-weighted Poisson regression for awareness and vaccination outcomes.

\begin{table}[htbp]
\caption{Adjusted prevalence ratios (PR) from multivariable Poisson regression (N~=~3,899)}
\centering
\small
\begin{tabular}{|p{3.5cm}|c|c|c|c|}
\hline
 & \multicolumn{2}{c|}{\textbf{Awareness}} & \multicolumn{2}{c|}{\textbf{Vaccination}} \\
\hline
\textbf{Variable} & \textbf{PR} & \textbf{95\% CI} & \textbf{PR} & \textbf{95\% CI} \\
\hline
\multicolumn{5}{|l|}{\textit{Rural residence (ref: Urban)}} \\
\hline
Yes & 0.933 & 0.867--1.005 & 0.780 & 0.547--1.112 \\
\hline
\multicolumn{5}{|l|}{\textit{Age group (ref: 25--29)}} \\
\hline
15--19 & 0.626*** & 0.561--0.699 & 0.514* & 0.305--0.865 \\
20--24 & 0.918* & 0.855--0.986 & 1.150 & 0.796--1.660 \\
\hline
\multicolumn{5}{|l|}{\textit{Education (ref: Upper secondary+)}} \\
\hline
Primary or less & 0.475*** & 0.330--0.683 & 0.327 & 0.060--1.770 \\
Lower secondary & 0.836*** & 0.754--0.927 & 0.422** & 0.227--0.785 \\
\hline
\multicolumn{5}{|l|}{\textit{Wealth quintile (ref: Q5 richest)}} \\
\hline
Q1 (poorest) & 0.576*** & 0.480--0.691 & 0.118*** & 0.039--0.358 \\
Q2 & 0.802*** & 0.726--0.885 & 0.585* & 0.365--0.936 \\
Q3 & 0.812*** & 0.743--0.888 & 0.437*** & 0.278--0.688 \\
Q4 & 0.890** & 0.829--0.956 & 0.631* & 0.421--0.945 \\
\hline
\multicolumn{5}{|l|}{\textit{Mass media exposure (ref: No)}} \\
\hline
Yes & 1.147 & 0.981--1.340 & 1.960 & 0.763--5.034 \\
\hline
\multicolumn{5}{|l|}{\textit{Ethnic minority (ref: No)}} \\
\hline
Yes & 0.860* & 0.745--0.992 & 0.743 & 0.307--1.801 \\
\hline
\multicolumn{5}{|l|}{\textit{Decision autonomy (ref: Low)}} \\
\hline
High & 1.234** & 1.066--1.429 & 1.614 & 0.861--3.025 \\
\hline
\end{tabular}
\label{tab:poisson}

\vspace{0.3cm}
\footnotesize{*$p<0.05$, **$p<0.01$, ***$p<0.001$. PR = prevalence ratio. N = 3,899 after listwise deletion.}
\end{table}

\textbf{Awareness.} After adjusting for all covariates, women in the poorest quintile had 42\% lower awareness than the richest (PR~=~0.576, 95\% CI: 0.480--0.691). Primary education was associated with a 53\% reduction (PR~=~0.475, 95\% CI: 0.330--0.683). Ethnic minority status (PR~=~0.860) and younger age (15--19: PR~=~0.626) were independently associated with lower awareness. High decision autonomy (PR~=~1.234) was protective.

\textbf{Vaccination.} The socioeconomic gradient was substantially steeper. Women in the poorest quintile had an 88\% lower prevalence of vaccination than the richest (PR~=~0.118, 95\% CI: 0.039--0.358). Women aged 15--19 had a 49\% lower prevalence than those aged 25--29 (PR~=~0.514). Lower secondary education showed a pronounced effect (PR~=~0.422).

\section{Erreygers decomposition analysis}

\subsection{Decomposition of awareness inequality}

Table~\ref{tab:decomp_awareness} presents the Erreygers decomposition results for HPV awareness inequality.

\begin{table}[htbp]
\caption{Erreygers decomposition of socioeconomic inequality in HPV awareness (Erreygers CI = 0.371)}
\centering
\small
\begin{tabular}{|p{3.5cm}|c|c|c|c|}
\hline
\textbf{Variable} & \textbf{dy/dx} & \textbf{CI$_x$} & \textbf{Contribution} & \textbf{\% of total} \\
\hline
Q1 poorest & $-$0.302 & $-$0.621 & 0.1440 & \textbf{38.8} \\
Q2 & $-$0.187 & $-$0.342 & 0.0544 & \textbf{14.6} \\
Rural & $-$0.050 & $-$0.412 & 0.0514 & \textbf{13.8} \\
Mass media & 0.063 & 0.180 & 0.0413 & \textbf{11.1} \\
Decision autonomy & 0.102 & 0.089 & 0.0327 & 8.8 \\
Never married & 0.062 & 0.237 & 0.0296 & 8.0 \\
Lower secondary & $-$0.102 & $-$0.317 & 0.0285 & 7.7 \\
Ethnic minority & $-$0.065 & $-$0.413 & 0.0161 & 4.3 \\
\hline
\end{tabular}
\label{tab:decomp_awareness}
\end{table}

The Erreygers-corrected CI for awareness was 0.371, confirming pro-rich concentration. Household wealth (Q1 poorest) was the dominant structural determinant, generating \textbf{38.8\%} of the total inequality. Geographic marginalization (rural residence) drove 13.8\%, lack of mass media access drove 11.1\%, and educational deficits contributed 7.7\% of the disparity.

\begin{figure}[htbp]
\centering
\includegraphics[width=0.90\textwidth]{Figures/fig8_decomposition_awareness.png}
\caption{Top 10 contributors to socioeconomic inequality in HPV awareness (Erreygers decomposition). Wealth status, rural residence, mass media exposure, and education are the dominant drivers.}
\label{fig:decomp_awareness}
\end{figure}

\subsection{Decomposition of vaccination inequality}

\begin{table}[htbp]
\caption{Erreygers decomposition of socioeconomic inequality in HPV vaccination (Erreygers CI = 0.116)}
\centering
\small
\begin{tabular}{|p{3.5cm}|c|c|c|c|}
\hline
\textbf{Variable} & \textbf{dy/dx} & \textbf{CI$_x$} & \textbf{Contribution} & \textbf{\% of total} \\
\hline
Q1 poorest & $-$0.125 & $-$0.621 & 0.0598 & \textbf{51.4} \\
Mass media & 0.042 & 0.180 & 0.0276 & \textbf{23.8} \\
Rural & $-$0.019 & $-$0.412 & 0.0199 & \textbf{17.1} \\
Lower secondary & $-$0.051 & $-$0.317 & 0.0143 & \textbf{12.3} \\
Q2 & $-$0.045 & $-$0.342 & 0.0129 & 11.1 \\
Sexual experience & $-$0.025 & $-$0.223 & 0.0119 & 10.2 \\
Decision autonomy & 0.033 & 0.089 & 0.0105 & 9.0 \\
\hline
\end{tabular}
\label{tab:decomp_vaccination}
\end{table}

The poorest wealth quintile alone is responsible for generating \textbf{51.4\%} of the total inequality in vaccine uptake. This confirms that financial insolvency overrides almost all other considerations at the point of care. Mass media exclusion generated 23.8\% of the disparity, while rural isolation (17.1\%) and educational deficits (12.3\%) acted as secondary structural barriers.

\begin{figure}[htbp]
\centering
\includegraphics[width=0.90\textwidth]{Figures/fig9_decomposition_vaccination.png}
\caption{Top 10 contributors to socioeconomic inequality in HPV vaccination (Erreygers decomposition). Wealth status alone explains over half of the observed inequality.}
\label{fig:decomp_vaccination}
\end{figure}

\subsection{Decomposition of awareness--vaccination gap inequality}

The Erreygers CI for the gap was $-$0.123, indicating pro-poor concentration: among aware women, those from lower socioeconomic strata were more likely to remain unvaccinated. The largest contributors were Q1 poorest (57.6\%), mass media (35.1\%), rural residence (16.4\%), and lower secondary education (14.4\%).

\begin{figure}[htbp]
\centering
\includegraphics[width=0.90\textwidth]{Figures/fig2_gap_by_ses.png}
\caption{Awareness--vaccination gap by wealth quintile among aware women. The gap is widest among the poorest, reflecting barriers beyond awareness.}
\label{fig:gap_ses}
\end{figure}

Figure~\ref{fig:gap_ses} visualizes this critical failure point: among women who already possess adequate knowledge of the HPV vaccine, those in the poorest quintile face the widest gap between intent and action. This confirms that informational campaigns alone are insufficient; without concurrent financial subsidies or inclusion in the national Expanded Programme on Immunization (EPI), the poorest women remain structurally locked out of cervical cancer prevention.